\section{Contexto Geral do Trabalho}
\label{sec:contextoGeral}

O trabalho partiu de um problema operacional concreto no Banco de 
Especialistas: as solicitações de alteração cadastral e de transição 
de categoria eram feitas por e-mail e executadas manualmente por um 
técnico, por meio de scripts SQL. O diagnóstico apresentado na Seção~\ref{sec:fluxo-antigo} 
evidenciou quatro consequências desse fluxo — risco de inconsistência, 
ausência de rastreabilidade, ineficiência operacional e dependência da 
equipe de TI —, reforçadas pelo volume de chamados: 141 registros 
entre maio de 2025 e agosto de 2026, dos quais 48,2\% correspondiam a 
ajuste cadastral ou de pontuação, e a categoria residual “Outro”, que 
incluía pedidos de mudança de categoria, respondia por 41,8\% do total.

Diante disso, o objetivo geral foi desenvolver um módulo full-stack de 
solicitações para o gerenciamento das operações de alteração cadastral 
e de classificação dos especialistas, com mecanismos de rastreabilidade 
e registro das operações, conforme os objetivos definidos na Seção~\ref{sec:objetivos}, 
preservando o comportamento das funcionalidades existentes por meio da 
estratégia de Additive Change na Seção~\ref{sec:additive-change}. Os objetivos 
específicos correspondentes — levantamento de requisitos e regras de negócio, 
modelagem da persistência, implementação do back-end, desenvolvimento das 
interfaces web, implementação da rastreabilidade, execução de testes e 
configuração da integração contínua — foram endereçados, respectivamente, nas 
Seções~\ref{sec:regras-negocio}, \ref{sec:modulo-solicitacoes}, \ref{sec:snapshot-auditoria}, \ref{sec:interfaces-web}, \ref{sec:testes-validacao} e \ref{sec:fluxo-antigo}.

O escopo foi deliberadamente delimitado na Seção~\ref{sec:escopo-limitacoes}: não 
contemplou a redefinição dos critérios de pontuação, a modificação do processo 
seletivo, a substituição do mecanismo de avaliação existente, a implementação de 
notificações por e-mail ou o tratamento de múltiplas solicitações simultâneas para 
um mesmo especialista. Essas exclusões delimitam o que o módulo se propôs a 
resolver — formalizar, controlar e rastrear as solicitações de transição de 
categoria e de complemento cadastral, com histórico das decisões e reversão 
controlada — e são a base sob a qual os resultados do Capítulo 3 devem ser 
avaliados.